\section{不可约性判别法}\label{sec:irreducibility-criteria}

为了在代数几何中有大量具体例子可用, 掌握研究特定多项式的方法是有益的. 例如, 如果我们考虑某个特定多项式, 如
\[
  x^{2}+y^{2}+z^{2}-1\in\mathbb{R}[x,y,z],
\]
那么能够迅速判断它是否不可约, 可能很有帮助. 在本章最后一节, 我们讨论两种判别给定多项式是否不可约的方法. 第一个结果可由前面的论述直接得出, 证明留给 Exercise~\ref{ex:irreducible-criteria-equivalence}.

\begin{Proposition}{不可约多项式的刻画}\label{prop:irreducible-polynomial-characterization}
  设 $f\in K[x_{1},\ldots,x_{n}]$. 下列条件等价.
  \begin{enumerate}[label=\arabic*.]
    \item $f$ 不可约;
    \item $f$ 是素元;
    \item $\langle f\rangle$ 是素理想;
    \item $K[x_{1},\ldots,x_{n}]/\langle f\rangle$ 是整环.
  \end{enumerate}
\end{Proposition}

\begin{Example}{{再论 \textmd{$y-x^{2}$} 在 \textmd{$K[x,y]$} 中不可约}}\label{ex:parabola-polynomial-irreducible-revisited}
  我们已经直接证明了 $y-x^{2}$ 不可约. 根据 Proposition~\ref{prop:irreducible-polynomial-characterization}, 由 $K[x]$ 是整环以及
  \[
    \frac{K[x,y]}{\langle y-x^{2}\rangle}\cong K[x].
  \]
  也能看出这一点.
\end{Example}

在实际应用中, Proposition~\ref{prop:irreducible-polynomial-characterization} 并非总有帮助. 证明一个理想是素理想或一个商环是整环, 通常和直接证明多项式不可约一样困难. 这个命题转换了问题, 却不一定简化了它. 下面的判别法称为 Eisenstein 判别法. 虽然表述稍复杂, 它在实际应用中却有用得多, 后面的例子将说明这一点.

\begin{Proposition}{Eisenstein 判别法}\label{prop:eisenstein-criterion}
  设 $R$ 是整环, 且 $f=a_{n}x^{n}+\cdots+a_{1}x+a_{0}\in R[x]$ 是满足以下条件的多项式.
  \begin{enumerate}[label=\arabic*.]
    \item 不存在非可逆元 $b\in R$ 使 $b\mid f$.
    \item 存在素元 $p\in R$ 使得
      \begin{itemize}
        \item 当 $i<n$ 时, $p\mid a_{i}$, 且
        \item $p^{2}\nmid a_{0}$.
      \end{itemize}
  \end{enumerate}
  则 $f$ 不可约.
\end{Proposition}

与 Proposition~\ref{prop:irreducible-polynomial-characterization} 不同, Eisenstein 判别法不是一个充要条件. 特别地, Eisenstein 判别法不能用来判别一元多项式 $f\in K[x]$ 是否不可约, 原因只是域中没有素元. 在证明 Eisenstein 判别法之前, 我们先给出几个应用的例子, 说明怎样将它用于域上的多元多项式. 习题中还有更多例子.

\begin{Example}{{\textmd{$f=x^{2}+y^{2}-1$} 在 \textmd{$\mathbb{R}[x,y]$} 中不可约}}\label{ex:eisenstein-unit-circle}
  若把 $f$ 看作 $R[y]$ 的元素, 其中 $R=\mathbb{R}[x]$, 则
  \[
    f=a_{2}y^{2}+a_{1}y+a_{0}
  \]
  其中 $a_{2}=1$, $a_{1}=0$, 且 $a_{0}=x^{2}-1=(x-1)(x+1)$. 注意, 这些系数在 $\mathbb{R}[x]$ 中没有公共的非常数因子, 因而 $f$ 满足 Eisenstein 判别法的第一个条件. 由于商环
  \[
    \frac{\mathbb{R}[x]}{\langle x-1\rangle}\cong\mathbb{R},
  \]
  \BookPairedNote{是整环, $p=x-1$ 在 $\mathbb{R}[x]$ 中是素元, 且满足 Eisenstein 判别法的第二个条件. 因此, $f$ 是不可约多项式.}%
    {这两个例子并不一定要在 $\mathbb{R}$ 上进行. 对于任何满足 $1\ne-1$ (即 $\operatorname{char}(K)\ne2$) 的域, 同样的论证都成立.}
\end{Example}

\begin{Example}{{\textmd{$f=x^{2}+y^{2}+z^{2}-1$} 在 \textmd{$\mathbb{R}[x,y,z]$} 中不可约}}
  \leavevmode\label{ex:eisenstein-unit-sphere}\strut\par\nopagebreak[4]\noindent
  和上一个例子一样, 写成
  \[
    f=a_{2}z^{2}+a_{1}z+a_{0}
  \]
  其中 $a_{2}=1$, $a_{1}=0$, 且 $a_{0}=x^{2}+y^{2}-1$. 这些系数在 $\mathbb{R}[x,y]$ 中没有公共的非常数因子, 所以 $f$ 满足 Eisentein 判别法的第一个条件. 令 $p=x^{2}+y^{2}-1$, 由上一个例子可知 $p$ 不可约, 且由于 $\mathbb{R}[x,y]$ 是 UFD, 因而 $p$ 是素元. 多项式 $p$ 满足 Eisenstein 判别法的第二个条件, 所以 $f$ 不可约.
\end{Example}

\begin{Example}{{\textmd{$f=x^{2}y^{2}+yz^{2}+x^{3}z^{2}\in K[x,y,z]$} 不可约}}
  \leavevmode\label{ex:eisenstein-three-variable-polynomial}\strut\par\nopagebreak[4]\noindent
  写成
  \[
    f=a_{3}x^{3}+a_{2}x^{2}+a_{1}x+a_{0}\in R[x]
  \]
  其中 $R=K[y,z]$, 且 $a_{3}=z^{2}$, $a_{2}=y^{2}$, $a_{1}=0$, $a_{0}=yz^{2}$. 注意, 这些系数在 $K[y,z]$ 中没有公共的非常数因子, 因此 $f$ 满足 Eisenstein 判别法的第一个条件. 令 $p=y\in R$. 由于
  \[
    \frac{R}{\langle p\rangle}=\frac{K[y,z]}{\langle y\rangle}\cong K[z]
  \]
  是整环, 可知 $p=y$ 在 $R$ 中是素元. 因为 $p$ 满足 Eisenstein 判别法的第二个条件, 所以 $f$ 不可约.
\end{Example}

\begin{Proof}\unskip\textbf{ of Proposition~\ref{prop:eisenstein-criterion}}\quad
  设 $f$ 和 $p$ 如命题所述, 并假设对某些 $g,h\in R[x]$ 有 $f=gh$. 我们必须证明 $g$ 或 $h$ 是 $R$ 中的可逆元. 由第一个条件, $f$ 不可被 $R$ 中任何非可逆元整除, 所以只需证明 $g$ 或 $h$ 是 $R$ 的元素.

  为作反证, 假设 $g=b_{k}x^{k}+\cdots+b_{0}$ 和 $h=c_{\ell}x^{\ell}+\cdots+c_{0}$ 的次数都为正. 考虑环同态
  \[
    \begin{aligned}
      \phi:R[x]&\to(R/\langle p\rangle)[x]\\
      \sum_{i\ge0}r_{i}x^{i}&\mapsto\sum_{i\ge0}[r_{i}]x^{i}.
    \end{aligned}
  \]
  根据 Proposition~\ref{prop:eisenstein-criterion} 对 $f$ 的假设, 可得
  \[
    \phi(f)=[a_{n}]x^{n}=\phi(g)\phi(h),
  \]
  且 $[a_{n}]\ne0$. 由于 $R/\langle p\rangle$ 是整环, $\phi(g)$ 和 $\phi(h)$ 必各由一个非零项组成 (Exercise~\ref{ex:single-term-product}), 并且由次数的可加性, 得 $\phi(g)=[b_{k}]x^{k}$, $\phi(h)=[c_{\ell}]x^{\ell}$. 特别地, 这意味着 $p\mid b_{0}$ 且 $p\mid c_{0}$, 所以 $p^{2}\mid b_{0}c_{0}=a_{0}$, 与对 $p$ 的假设矛盾.
\end{Proof}

\BookExerciseHeading
\begin{BookExercises}
  \item\label{ex:irreducible-criteria-equivalence} 引用前几节的相关结果, 证明 Proposition~\ref{prop:irreducible-polynomial-characterization}.
  \item\label{ex:single-term-product} 设 $R$ 是整环, 且 $g,h\in R[x]$ 的乘积 $gh$ 只有一个非零项, 即 $gh=ax^{n}$, 其中 $a\in R$ 非零且 $n\in\mathbb{N}$. 证明 $g$ 和 $h$ 都只有一个非零项.
  \item\label{ex:determinant-two-by-two-irreducible} 证明 $wx-yz\in K[w,x,y,z]$ 不可约.
  \item\label{ex:mixed-cubic-polynomial-irreducible} 证明 $xyz+x^{2}z^{2}+yz^{3}+x\in K[x,y,z]$ 不可约.
  \item\label{ex:quadratic-sphere-any-dimension-irreducible} 假设 $\operatorname{char}(K)\ne2$ 且 $n\ge2$. 证明 $x_{1}^{2}+\cdots+x_{n}^{2}-1$ 在 $K[x_{1},\ldots,x_{n}]$ 中不可约.
  \item\label{ex:fermat-polynomial-irreducible} 假设 $n\ge3$ 且 $m\ge1$. 证明 $x_{1}^{m}+\cdots+x_{n}^{m}\in\mathbb{C}[x_{1},\ldots,x_{n}]$ 不可约.
\end{BookExercises}
